\chapter{Agrupamento e Redução de Dimensionalidade}

\section{Métodos não supervisionados}

Aprendizado não supervisionado encontra estrutura nos dados sem uma
variável alvo. O \hayashi{} fornece algoritmos de agrupamento,
embutimentos não lineares e ferramentas baseadas em PCA.

\section{K-Means}

K-Means particiona observações em $K$ grupos esféricos. É rápido e
escala bem, mas assume grupos de tamanho e forma similares.

\begin{lstlisting}[caption={K-Means}]
set_seed(42)
let n = 300
generate df x1 = rnormal(n, 0, 0.7)
generate df x2 = rnormal(n, 0, 0.7)

let m_km = kmeans(df, x="x1,x2", k=3)
print(m_km)
\end{lstlisting}

Use `m_km.labels` para ler o cluster atribuído e `m_km.centroids` para
inspecionar as médias dos clusters.

\section{Agrupamento hierárquico e DBSCAN}

Agrupamento hierárquico constrói uma árvore de grupos aninhados.
`hclust` pode ser cortado em qualquer número de grupos.

\begin{lstlisting}[caption={Agrupamento hierárquico}]
let m_hc = hclust(df, x="x1,x2", k=3, linkage="ward")
print(m_hc)
\end{lstlisting}

DBSCAN encontra regiões densas e rotula observações esparsas como ruído.

\begin{lstlisting}[caption={DBSCAN}]
let m_db = dbscan(df, x="x1,x2", eps=0.8, min=5)
print(m_db)
\end{lstlisting}

\section{Modelos de Mistura Gaussiana}

GMM modela cada cluster como uma gaussiana multivariada. É uma variante
suave do K-Means que permite clusters elípticos.

\begin{lstlisting}[caption={Agrupamento por mistura gaussiana}]
let m_gmm = gmm_clust(df, x="x1,x2", k=3)
print(m_gmm)
\end{lstlisting}

\section{Embeddings não lineares}

t-SNE e UMAP reduzem dados de alta dimensão para duas ou três
dimensões preservando a estrutura local. São úteis para visualização.

\begin{lstlisting}[caption={t-SNE}]
let m_tsne = tsne(df, x="x1,x2,x3", k=2, perplexity=30)
print(m_tsne)
\end{lstlisting}

\begin{lstlisting}[caption={UMAP}]
let m_umap = umap(df, x="x1,x2,x3", k=2, neighbors=15)
print(m_umap)
\end{lstlisting}

\section{Agrupamento espectral}

Agrupamento espectral usa a estrutura de autovalores de um grafo de
similaridade. Pode capturar formas de cluster não convexas onde
K-Means falha.

\begin{lstlisting}[caption={Agrupamento espectral}]
let m_sc = spectral(df, x="x1,x2", k=3)
print(m_sc)
\end{lstlisting}

\section{PCA, fator e biplot}

PCA e análise fatorial reduzem a dimensionalidade linearmente. `biplot`
visualiza observações e variáveis no mesmo espaço de PCA.

\begin{lstlisting}[caption={PCA}]
let m_pca = pca(df, x="x1,x2,x3", n=2)
print(m_pca)
\end{lstlisting}

\begin{lstlisting}[caption={Biplot}]
let m_biplot = biplot(df, x="x1,x2,x3", type="symmetric")
print(m_biplot)
\end{lstlisting}

\section{Estimativa de densidade kernel}

`kde` estima a densidade de uma única variável com um kernel escolhido.

\begin{lstlisting}[caption={KDE}]
let m_kde = kde(df, x1, bw=0.5, kernel="gaussian")
print(m_kde)
\end{lstlisting}
